\chapter{Appendix: Supplementary Material}
\label{appendix:supplementary}


\section{The wider substrate ladder}
\label{app:wider-ladder}
\bigskip

Section~\ref{sec:act1-assignment} rests the weight-assignment claim on the
weight-permuted control and reports two further rungs as corroboration. Their
values are here. The clustering-preserving and modularity-preserving rewires
each hold fixed everything the degree-preserving rewire does and one further
property of the topology besides; 
the random gaussian rung is included for completeness and takes no part in any
comparison, since it is not a rung of the criticality-matched ladder and matches
the connectome's edge count only in expectation.

Holding progressively more of the connectome's topology fixed recovers none
of what the assignment supplies. At $N = 448$ the strongest null of all six is the
clustering-preserving rewire at
2.081, % TIER0 3.1(b)
above every rung the figures draw, and the connectome's margin over it is
1.479-fold; % TIER0 3.1(b)
at $N = 1000$ Erdős–Rényi remains the strongest and the margin is
1.635-fold. % TIER0 3.1(b)
The separation therefore survives at both parcellations against every null
available, but its size depends on which null set it is measured against, and it
is quoted here with that set named. On $|\lambda_1|$ all three additional rungs
sit inside the null band at both scales, so the reading of
Chapter~\ref{chapter:act-I} is unchanged: a large Perron root over a bulk close
to everyone's.

These three rungs are not part of the four-substrate ladder. They appear in no
figure of Chapter~\ref{chapter:act-I}, the published spreads in
Table~\ref{tab:act1-ladder} are the ladder's and are not recomputed across the
wider set, and the near-identity of the absolute bulks stated at $N = 448$ is a
statement about the four.


\begin{table}[ht]
    \centering
    \caption[The three off-ladder substrates]{\textbf{The three substrates
    outside the four-rung ladder.} Aggregation and units as
    Table~\ref{tab:act1-ladder}: medians over ten seeds, raw units of $\mathbf{W}$,
    absolute bulk the product of the medians. All three are resampled at every
    seed, so unlike the connectome the two aggregation orders differ on every
    row, and only the product of medians satisfies the gap-ratio identity.}
    \label{tab:act1-offladder}
    \begin{tabular}{lrrrr}
        \toprule
        substrate & absolute bulk & $|\lambda_1|$ & $\mathrm{bulk95}$ & gap ratio \\
        \midrule
        \multicolumn{5}{l}{\emph{$N = 448$}} \\
        random gaussian & 0.0593 & 0.1104 & 0.5374 & 1.861 \\ % TIER0 3.1(b)
        clustering rewire & 0.0582 & 0.1212 & 0.4806 & 2.081 \\ % TIER0 3.1(b)
        modularity rewire & 0.0598  & 0.1167 & 0.5123 & 1.952 \\ % TIER0 3.1(b)
        \midrule
        \multicolumn{5}{l}{\emph{$N = 1000$}} \\
        random gaussian & 0.0558 & 0.1357 & 0.4113 & 2.432 \\ % TIER0 3.1(b)
        clustering rewire & 0.0581 & 0.1377 & 0.4218 & 2.371 \\ % TIER0 3.1(b)
        modularity rewire & 0.0604 & 0.1353 & 0.4462 & 2.241 \\ % TIER0 3.1(b)
        \bottomrule
    \end{tabular}
\end{table}

% full substrate family adjacency matrices
\begin{sidewaysfigure}
    \centering
    \includegraphics[width=0.95\textheight]{S4.png}
    \captionsetup{font=small}
    \caption[The full substrate family, as graphs]{%
    \textbf{The full substrate family.} All seven substrates, drawn under the
    node ordering of Figure~\ref{fig:F19}: hemisphere, then community, with
    communities detected once on the connectome and reused throughout.
    \textbf{Top row} binary adjacency with the degree marginal above each panel;
    \textbf{bottom row} weighted adjacency, all seven on one logarithmic colour
    scale. Columns run in the order of
    Table~\ref{tab:methods-preservation}, most preserved first. The
    criticality-matched ladder Chapter~\ref{chapter:act-I} draws is the
    connectome, the weight-permuted control, the degree-preserving rewire and
    Erdős–Rényi, in columns one, two, five and six; the clustering-preserving
    and modularity-preserving rewires in columns three and four each hold fixed
    everything the degree rewire does and one further property of the topology
    besides, and the random gaussian rung in column seven fixes only the node
    count and the expected density, so it is the one substrate whose edge count
    varies between seeds. Every randomised substrate takes its weights from the
    connectome's own multiset, so all seven are non-negative. This figure shows
    what the substrates look like and makes no comparison among them: the ladder
    statistics of Table~\ref{tab:act1-ladder} are the ladder's and are not
    recomputed across this set, and the spectra of the three additional
    substrates are in Table~\ref{tab:act1-offladder}. Randomised substrates are
    drawn at the seed used throughout this chapter.}
    \label{fig:S4}
\end{sidewaysfigure}





\section{A pre-registered test of one mechanism}
\label{app:mechanism}
\bigskip

Chapter~\ref{chapter:act-I} localises the spectral gap to weight assignment and
claims no mechanism for it. One candidate mechanism was specific enough to
test, and was tested. The prediction, its falsifier and the decision rule were
committed before any of the analysis ran.

The candidate is the weighted rich club. For a symmetric non-negative matrix
the Perron root is the maximum of the Rayleigh quotient over non-negative
vectors, so heavy weights on edges sharing endpoints reinforce one another in
it while scattered heavy weights do not. If the connectome's heavy edges ran
between its hubs they would share endpoints, and a weight permutation, which
holds the graph and the multiset fixed, would destroy that arrangement and the
Perron root with it.

Three predictions were registered. Edge weight would correlate positively with
endpoint degree product in the connectome and at zero in the control; the
heaviest edges would carry a larger fraction of the Perron root in the
connectome than in the control; and the ratio of maximum node strength between
them would be smaller than their Perron ratio, excluding a single dominant node
as the explanation.

The first fails, in the direction opposite to the prediction, and
Figure~\ref{fig:S3} is the measurement. The second was posed on a fraction of
each substrate's own Perron root, which divides by the one quantity these
substrates differ in, so it cannot be read: the same normalisation this thesis
retracts in Section~\ref{sec:act1-assignment}. The third fails and
leaves an account standing rather than excluding one. Weight assignment
concentrates node strength: the connectome's maximum node strength is
2.54 times % TIER0 3.14
the control's while its Perron root is 1.65 times, % TIER0 3.14
on identical mean strength. That is a candidate this thesis does not test. It
was not the registered hypothesis, it arrives from a measurement designed to
exclude it rather than to examine it, and adopting it here would be the
post-hoc move the registration exists to prevent.

The registered decision rule returned equivocal, so nothing in the claims
register changed and the limit stated in
Section~\ref{sec:act1-assignment} stands as written.

\bigskip
\begin{figure}[ht]
    \centering
    \includegraphics[width=\textwidth]{S3.png}
    \caption[The connectome's heavy edges avoid its hubs]{%
    \textbf{The connectome's heavy edges run between its lower-degree parcels.}
    Mean edge weight against the product of the two endpoints' binary degrees,
    in twenty equal-count bins cut on the connectome's products so that both
    substrates bin identically. The weight-permuted control holds the graph and
    the weight multiset fixed, so any trend it shows is sampling noise, and it
    shows none: Spearman
    $-0.0033$, % TIER0 3.14
    with a bootstrap interval covering zero. The connectome's trend is negative,
    Spearman $-0.1242$ % TIER0 3.14
    with an interval excluding zero. A weighted rich club would put the trend
    the other way. 
    Ten thousand edge resamples, 95\% percentile intervals, $N = 448$. % TIER0 3.14
    }
    \label{fig:S3}
\end{figure}




\section{Sign selects the basis}
\label{app:sign-basis}
\bigskip
% Carries A2.3, the one claim in the act that is deliberately NOT a
% contribution. Confirmatory of Krauss 2019. Moved here from chapter 5 as
% secondary material bounding a claim made there; F5 moves with it.
% Sources: fact sheet 11, plus audit item 23 for the nested coverage.
% FORBIDDEN here (sheet 11): "the manifold lives in low-frequency graph
% harmonics"; any ladder reading of this capture; "sign-gating of the manifold
% transition is a new result"; "balancing the signs improves the
% representation"; quoting the capture values as evidence of how much structure
% the bases hold; treating this as a result about the sign fraction f, which is
% chapter 6's intervention; "compact bulk"; any graded curvature account.

Section~\ref{sec:act2-perron} found the low-frequency harmonics of the graph
Laplacian standing well above chance among the connectome's fluctuations while
the dominant eigenvectors of $\mathbf{W}$ fell below it. That is easy to read as a fact
about the connectome, and it is not. It is a fact about non-negativity, and the
way to show as much is to remove the non-negativity and leave everything else
where it is.

With non-negative weights the leading eigenmodes of $\mathbf{W}$ have been absorbed into
the rank-one term of Equation~\ref{eq:gram-split}, which leaves the fluctuations
to be organised by something other than the dynamical modes themselves.
Balancing the signs of the weights destroys the Perron guarantee, so there is no
dominant common mode for those eigenmodes to be absorbed into, and the
prediction is simply that they should come back.

They do. Table~\ref{tab:app-basis} repeats the basis comparison across three
weight conditions on one substrate: the empirical weights used throughout
Chapter~\ref{chapter:act-II}, the same magnitudes with their signs balanced, and
gaussian weights. The ordering of the two structural bases swaps between the
first condition and the other two, it swaps in the same direction on all three
tasks, and the topology never changes: the same graph, the same edges, and in
the first two conditions the same weight magnitudes, with only the signs
redistributed. Figure~\ref{fig:F5} draws the clearest of the three cases across
the whole range of $k$.

The swap is the measurement, and the magnitudes are not. On the all-positive
substrate neither basis holds much of anything, the whole range across bases and
tasks running
0.04 to 0.17, % TIER0 3.12
so the fluctuations are not being shown to live in graph harmonics. Which basis
leads is what the comparison establishes; how much either one captures is a
separate question, and the answer to that one is not much.

That low capture is what the other probe predicts. Supercritically the readout
has around 413 % TIER0 3.6
of 448 directions available to it on this substrate, so the fluctuations are
spread over hundreds of directions and no ten-vector basis should be able to
hold them. The agreement runs to the other end as well: where a condition's
effective rank is low the capture is high, reaching
0.886 % TIER0 3.12
for the gaussian weights on NARMA-10. Two probes measuring different things on
different bases agree about which conditions spread their activity widely, and
neither was fitted to the other.

Two limits. The first is priority. That weight sign composition gates which
basis the fluctuations occupy is largely pre-empted by Krauss et
al.~\cite{krauss2019weight}, and it is reported here as confirmatory; the
counting argument of Section~\ref{sec:act2-counting} needs the basis established
first. The second is scope. The signed and gaussian conditions
are weight families compared side by side. Nothing has been done to the
connectome to produce them, and the question of what happens when negative
weights are introduced as an intervention on the real substrate belongs to
Chapter~\ref{chapter:act-III}, where it is used to remove the common mode and to
check whether its consequences go with it.

\begin{table}[ht]
    \centering
    \caption[Which structural basis the fluctuations occupy]{\textbf{Which
    structural basis the fluctuations occupy, across three weight conditions and
    three tasks.} Fraction of time-centred state variance held by the leading
    ten vectors of each basis, connectome substrate,
    $N = 448$, % TIER0 2.1
    medians over ten seeds. Each condition is measured at its own supercritical
    operating point, 3.0526, 2.5263 and 1.2632 % A2 2.3
    respectively, since the three do not become supercritical at the same
    nominal radius; the three columns are therefore not one $\sigma$. The harmonics lead in every cell of the first condition and the $\mathbf{W}$ eigenmodes
    lead in every cell of the other two. The magnitudes are small: on the
    all-positive substrate neither basis holds much of anything, and the
    ordering is what the table establishes. Recomputation against the published table
    agrees to 0.00045 % A2 2.3
    across all eighteen cells.}
    \label{tab:app-basis}
    \begin{tabular}{llcc}
        \toprule
        & & \multicolumn{2}{c}{fraction held at $k = 10$} \\
        \cmidrule(lr){3-4}
        weight condition & task & harmonics & $\mathbf{W}$ eigenmodes \\
        \midrule
        empirical (\textbf{non-negative}) & memory capacity & \textbf{0.056} & 0.009 \\ % TIER0 3.12
        & NARMA-10 & \textbf{0.040} & 0.002 \\ % TIER0 3.12
        & Lorenz & \textbf{0.168} & 0.004 \\ % TIER0 3.12
        \addlinespace
        empirical (\textbf{signs balanced}) & memory capacity & 0.017 & \textbf{0.060} \\ % TIER0 3.12
        & NARMA-10 & 0.005 & \textbf{0.527} \\ % TIER0 3.12
        & Lorenz & 0.019 & \textbf{0.081} \\ % TIER0 3.12
        \addlinespace
        gaussian (\textbf{signs balanced}) & memory capacity & 0.020 & \textbf{0.040} \\ % TIER0 3.12
        & NARMA-10 & 0.006 & \textbf{0.886} \\ % TIER0 3.12
        & Lorenz & 0.012 & \textbf{0.522} \\ % TIER0 3.12
        \bottomrule
    \end{tabular}
\end{table}

% Figure F5.
\begin{figure}[ht]
    \centering
    \includegraphics[width=\textwidth]{F5.png}
    \caption[Sign composition selects which structural basis the fluctuations
    occupy]{%
    \textbf{Which structural basis the fluctuations occupy is set by the weight
    signs, not by the topology.} Fraction of time-centred state variance held by
    the leading $k$ vectors of each basis, connectome substrate, Lorenz task,
    $N = 448$, % TIER0 2.1
    medians over ten seeds, each panel at its own condition's supercritical
    operating point. \textbf{(a)} With \textbf{non-negative} weights the graph
    harmonics lead throughout and the dominant $\mathbf{W}$ eigenmodes run below chance.
    \textbf{(b)} With the \textbf{same weight magnitudes} and balanced signs the
    ordering has already reversed. \textbf{(c)} With \textbf{gaussian} weights
    the eigenmodes dominate outright. The rule marks $k = 10$, where the values
    of Table~\ref{tab:app-basis} are read; the swap is visible across the whole
    curve, and the topology is identical in all three panels. On the all-positive substrate neither basis holds much in
    absolute terms; the ordering is what the panels establish. The degree-preserving rewire, the only
    other substrate this capture holds, shows the same swap at the same three
    operating points, reading
    0.168 against 0.002, % A2 F5
    0.011 against 0.102, % A2 F5
    and 0.011 against 0.355; % A2 F5
    two substrates cannot speak to the null ladder.}
    \label{fig:F5}
\end{figure}




\section{Peak memory capacity against ridge strength}
\label{app:peak-parity}
\bigskip

Section~\ref{sec:act3-memory-sigma} states parity at the peak and reports the
interval behind it. The comparison is here in full. Each substrate is taken at
its own best spectral radius and differences are paired within seed, so every
comparison shares its input weights and its input series, and the whole panel is
repeated at five ridge strengths spanning five orders of magnitude. The five are
shown because usable dimensionality tracks measured memory at
$+0.999$ % TIER0 3.3
across all of them, on the seed-median statistic over
$n = 52$; % TIER0 3.4
raising the regularisation does not break that correspondence provided it is
raised in both places at once, so the readout's regularisation is not what
produces the result in either direction.

The smallest ridge is also where this thesis meets an ordering in the literature
that appears to contradict it. Nearest the pseudoinverse limit, peak memory
orders Erdős--Rényi, degree-matching, weight-permuted, connectome, which is the
spread-beats-compact ordering of Aceituno et al.~\cite{aceituno2020tailoring}
reproduced exactly on these substrates. Section~\ref{sec:act3-memory-sigma} reads
the two results as statements about different places on the operating range: spread wins at
the peak, and a large spectral gap wins across the range.

\begin{figure}[ht]
    \centering
    \includegraphics[width=\textwidth]{F10.png}
    \caption[At its own optimum the connectome's memory is at parity with the
    nulls]{\textbf{At its own optimum the connectome's memory is at parity with the
    nulls, not below them all.} Paired per-seed differences in peak memory capacity,
    each substrate taken at its own best spectral radius, ten seeds, with input
    weights and input series shared within a pair. Bars are $95\%$ $t$-intervals and
    filled markers are intervals excluding zero. \textbf{(a)} In memory-capacity
    units and \textbf{(b)} as a percentage of each null's own peak. The deficit is
    $2$ to $6\%$, and it excludes zero at all five ridge strengths against
    Erdős--Rényi and against the weight-permuted assignment control, but at only one
    of five against degree-matching, where the estimate also changes sign: the open
    markers. The five strengths span five orders of magnitude and are shown because
    usable dimensionality tracks measured memory at $+0.999$ across all of them
    (seed-median statistic, $n = 52$), so the readout's regularisation is not what
    produces the result.}
    \label{fig:F10}
\end{figure}



\section{The memory margin across network scale}
\label{app:scale}
\bigskip


Figure~\ref{fig:F9} reports a memory margin that holds across a 2.2-fold change in 
network size. Two things changed between the two runs, and they are separated by a 
control. The run length was scaled from
$3000$ to $6000$ % TIER0 2.5
so that the ratio of retained timesteps to units stayed near
$5.5$, % TIER0 2.5
and the ridge was reparameterised to scale with the Gram matrix rather than held
at an absolute value. At fixed $N$ the reparameterisation alone moves the
supercritical margin from
$4.35$ to $4.40$, % TIER0 2.4
so anything beyond that is attributable to size.
Two predictions stated before the run held: the margin does not vanish, and the
larger network does not escape the ceiling, peak usable dimensionality running at
$0.971$ for the connectome and $0.984$ to $0.999$ % TIER0 2.4
of its maximum for the nulls.

One thing this run was designed to settle and did not. It carried a
pre-registered falsification test, on the prediction that if the bulk radius is
what orders the ladder, then a reversal in that quantity between the two sizes
should reverse the memory ordering with it. The memory ordering did not move.
But the reversal the test rests on is itself not significant at the larger size
($p = 0.16$), % TIER0 2.4
so the test was flawed rather than passed: the correct classification is
inconclusive, and it is inconclusive because the predictor was noisy and not
because the outcome was. The registration asserted that the test was on the sign
rather than the size without first checking that the sign was established. The
question it was meant to answer is asked again in
Section~\ref{sec:act3-lesion}, by an intervention that removes the
proposed mechanism outright and therefore does not depend on a predictor moving.




\section{A second mode of generative failure}
\label{app:second-failure}
\bigskip
% Destination for chapter 6 section 6.3's zero-VPT paragraph, moved here as a
% negative result bounding A3P.4. The bounding statement stays in the main text
% beside the claim (CONVENTIONS working rule 8); the counts are here.
% NUMERALS GATED: every count below is from the read-only audit of 30 Aug as
% reproduced by verification V4 on 3 September, and is filed nowhere above rank
% 5. DO NOT SET until GAPS B26 is committed. Unit: the seed-cell,
% 1160 = 4 variants x 29 sigma x 10 seeds, three draws collapsed because at
% f = 0 the sign transform is the identity.
% The reason chapter 6 previously gave for separating the two modes is
% WITHDRAWN: the collapse flag IS persisted per seed and draw, and the two can
% be matched. See V4.
% FORBIDDEN: reading the two modes as one event; "50% of replicates".

Section~\ref{sec:act3-prediction-sigma} reports that the connectome alone avoids
the period-two collapse, and bounds that result with a second and cruder measure
on which the ordering reverses. Counting seeds that record a valid prediction
time of exactly zero at any operating point, the connectome has
five of ten, % GAPS B26
the weight-permuted control
three, % GAPS B26
Erdős--Rényi
two, % GAPS B26
and the degree-matched null none at all, its lowest value anywhere in the sweep
being
$0.068$. % GAPS B26

Three things bound how much should be read into that. The failures are confined
to the upper sweep, the lowest operating point at which any connectome seed
records a zero being
$\sigma = 3.6$, % GAPS B26
so none occurs at or below the point at which the near-critical comparison of
Section~\ref{sec:act3-prediction-sigma} is made. They are never a whole
condition failing: at worst half the seeds of a condition score zero while the
rest predict at between one and three Lyapunov times, in line with their
neighbours. And they are a different event from the collapse rather than a
milder form of it, which can now be shown rather than assumed. Both flags are
recorded per seed, so the forty seeds at zero sign fraction cross-tabulate
directly:

\begin{center}
\begin{tabular}{lcc}
    \toprule
    & \multicolumn{2}{c}{zero prediction time} \\
    \cmidrule(lr){2-3}
    collapses to period two & no & yes \\
    \midrule
    no  & $25$ & $6$ \\
    yes & $5$  & $4$ \\
    \bottomrule
\end{tabular}
\end{center}

All five of the connectome's zero-scoring seeds sit in the upper-left and
upper-right cells, since it has no collapsing seed at all; Erdős--Rényi's two
are both collapsing seeds. The two failure modes are therefore distinct, they
order the substrates differently, and both orderings are correct measurements of
different things.


\section{Regime gating against graded geometry}
\label{app:regime-gating}
\bigskip

Section~\ref{sec:act3-branch} states that predictive capacity is gated by which
dynamical regime a network occupies rather than graded by how curved its
trajectory is. The comparison behind that statement is here.

Across the full sweep in operating point and sign fraction, a single binary
variable recording which curvature spike a cell belongs to explains
$R^2 = 0.364$ % TIER0 3.10
of the variance in valid prediction time. The full continuous curvature measure
explains
$0.370$. % TIER0 3.10
The entire range of geometries lying between the two regimes is therefore worth
less than one percentage point beyond the bit, and the bit is what the
relationship consists of.

The graded account it replaces is not a straw man but an earlier reading of these
same data. Pooled across all cells, curvature and prediction correlate at
$-0.78$, % TIER0 3.10
which reads as a relationship in which capacity tracks how straight the manifold
is. It is cluster mixing. Within the smooth cluster alone the residual
correlation is
$+0.145$, % TIER0 3.10
opposite in sign, and no binning of the data recovers a graded path between the
clusters. The pooled figure measures the separation between two groups and not a
trend within either.




\section{Attractor retention in long rollouts}
\label{app:long-rollout}
\bigskip

Valid prediction time is measured over short rollout windows and scores how long
a released network stays with the true trajectory before leaving it. A longer
question is whether a network left running on its own output continues to
produce the right kind of trajectory at all, which is measured over a rollout of
$3000$ % A3P 2.6
steps with the first
$500$ % A3P 2.6
discarded as settling, some
$67.9$ % A3P 2.6
Lyapunov times scored, and reported as a rate over cells rather than per cell
(Section~\ref{sec:methods-closed-loop}). This section reports that measurement.

A pre-registration written before the code existed made two claims about the
long rollout: that the connectome would retain the true attractor's statistics to
a higher sign fraction than the nulls, and that when substrates failed they would
fail by distorting the attractor's shape while preserving its scale.

The first is confirmed. Once a third of the weights are negative the connectome
holds a faithful climate in
$0.43$ % A3P 6.1
of cells against the nulls'
$0.14$, % A3P 6.1
about three times the rate (Figure~\ref{fig:F17}).

The second is refuted, and the refutation is the more informative result. The
ratio of the free run's spread to the true trajectory's spread is sharply
bimodal, with cells either collapsing below a twentieth of it or holding within a
tenth of unity and almost none between. A failed rollout is therefore a collapse
of scale and not a distortion of shape, which is the fixed-point branch of the
map of Section~\ref{sec:act3-branch} seen in the closed loop, on a different
quantity and a different capture from anything in that chapter.

\begin{figure}[htbp]
    \centering
    \includegraphics[width=\textwidth]{F17.png}
    \caption[Released to run alone, the connectome holds the true attractor
    further, and failure is a collapse of scale]{\textbf{Released to run alone,
    the connectome holds the true attractor further, and failure is a collapse of
    scale.} Free-running Lorenz rollout, $N = 448$, at $\sigma = 2$, against the
    sign fraction $f$. \textbf{(a, b)} Example rollouts either side of the
    separation. \textbf{(c)} The ratio of free-run spread to true spread, which
    refutes the second registered claim: of $440$ cells, $229$ fall below $0.05$
    and $158$ lie within a tenth of unity, with $11$ between, so failure is a
    loss of scale rather than a distorted shape. The registered description,
    wrong shape at right scale, fits $34$ cells. \textbf{(d)} The fraction of
    cells whose free run keeps the true attractor's statistics, confirming the
    first registered claim. \textbf{(e)} Where collapsed rollouts come to rest,
    most of them several z-scored units from the origin, against an attractor
    spanning about $\pm 2.5$ per coordinate. Panel (a) is read from the frozen
    capture and the rate in the text from a separate fresh capture; the two are
    not one population and are not combined.}
    \label{fig:F17}
\end{figure}



% ------------------------------

% S1
\begin{figure}[ht]
    \centering
    \includegraphics[width=\textwidth]{S1.png}
    \caption[Figure~\ref{fig:F1} rebuilt at the larger parcellation]{
    \textbf{The spectral gap at $N = 1000$.} As Figure~\ref{fig:F1} in every
    respect, same substrate family, same null ladder and same builder, at the
    $N = 1000$ % TIER0 2.1
    self-built consensus. \textbf{(a--d)} The four bulk bands again sit at the
    same place while $\lambda_1$ separates, and the connectome's gap ratio again
    stands at least 1.63-fold % TIER0 3.1(b)
    clear of every rung of the four-substrate ladder, 3.99 % TIER0 3.1
    against 2.30 to 2.44. % TIER0 3.1
    The rows are in ladder order and their gap ratios are \emph{not} monotone
    here, reading 3.99, 2.39, 2.30 and 2.44 % A1 S1 caption
    down the column, because the ordering of the three nulls by
    $\mathrm{bulk95}$ reverses between the two parcellations. What is
    scale-robust is the connectome's separation from the nulls, not the
    arrangement among them, and not the near-identity of the bulks, whose spread
    loosens from 4.4\% at $N = 448$ % TIER0 3.1
    to 6.4\% here, on the product of the medians. % TIER0 3.1
    \textbf{(e)} $\mathrm{bulk95}$ read off the 95\% crossing:
    0.251 % A1 S1 caption
    against 0.410 to 0.435. % A1 S1 caption
    \textbf{(f)} The per-seed gap ratio; the connectome's ten seeds again
    coincide exactly.}
    \label{fig:S1}
\end{figure}
